\documentclass[a4paper,10pt]{article}
\usepackage[utf8]{inputenc}
\RequirePackage{amsthm,amsmath,amssymb,amscd,xcolor}
\def\D{^{C}D_t^\alpha}
\def\W{\mathbb{W}}
\def\H{\mathbb{H}}
\def\K{\mathbb{K}}
\def\E{\mathbb{E}}
\def\cL{\mathcal{L}}
\def\R{\mathbb{R}}
\def\ud{\text{d}}
\newcommand{\Norm}[1]{\left|\left|  #1   \right|\right|}
\begin{document}

Here are comments for {\it 9cf2008f} on the file -- {\it thesis.tex}:

\begin{enumerate}
  \item Address point 1 of ``comment3\_d117e49.pdf".
  \item Address point 2 (b) with ``functions" of comment3.
  \item Address point 2 (c) with ``refers to" or you can see ``we use ... to denote the Fourier transform".
  \item Address point 2 (d) with $z\in \mathbb{Z}$.
  \item Eq. (2.10): $x\in\R$ should be $x\in\R^d$. Check all similar mistakes throughout the draft. For example, such
    typos appeared many times in Section 1.3.
  \item No need to add ``then" for point 4 of comment3. The usage is ``if..., then...", or ``when, ...".
  \item (2.9), forgot a period; check all similar problems.
  \item One line above (2.11): ``recall the subadditivity of the square root, namely, $\sqrt{a+b}\le \sqrt{a}+\sqrt{b}$
    \textcolor{magenta}{for all $a,b\ge 0$}".
  \item[] Accordingly, line -3 of the proof of Proposition 2.4: ``applying the identity ..." $\to$ ``thanks to the
    subadditivity of the square root".
  \item The second centered equations on p. 11 ($\mathbb{P}\left(\cdots\right)\le ...$), you may change them into one line.
\end{enumerate}

\end{document}
